\section{What the fits can, and cannot, be trusted for}
\label{sec:reliability}

\begin{investigation}{I8}{Do the delivered fits rank individual specimens
correctly, or only describe them correctly on average?}
\invQuestion{Against expert-corrected joint annotations, how many
joint-derived features reliably distinguish one specimen from another?}
\invMethod{Correlation of 49 joint-derived features against expert ground
truth; a parallel comparison of surface-landmark-derived traits against
rig-joint-derived traits at genus level.}
\invResult{Of 49 joint-derived features, only one reaches a correlation of
$0.5$ with expert ground truth. Registration of an independent
ethanol-preserved worker corpus sits at a measurable noise floor, at which
conspecific nest-mates are only $7\%$ closer than unrelated specimens; the
traits that fail the reliability threshold are precisely the distal leg and
antenna segment lengths that \F{F5} identified as under-sampled. On the same
fits, anatomical landmarks placed directly on the mesh surface carry more
genus-level signal than rig-derived joint distances in the earlier (provisional, \S\ref{sec:genus}) analysis, since a
rig joint is a rotation pivot rather than an anatomically defined point.}
\invVerdict{HOLD}{Population-level aggregates over many specimens are
supported by the current fits; per-specimen and comparative claims from
joint-derived measurements are not. For this reason the delivered
measurement layer reports surface-derived traits rather than rig lengths.}
\end{investigation}

A further consequence concerns the validation instrument itself, not only
the fitting pipeline: the project's own synthetic reliability estimate did
not survive contact with this expert ground truth. This is a second, more
consequential instance of the metric-versus-mechanism pattern in
\S\ref{sec:methodological-discipline}, here applied to the validation
instrument rather than to a fitting intervention, and it is the reason
this report treats expert-verified ground truth, wherever available, as the
higher-authority reference throughout.

\section{Validity is scale-dependent}
\label{sec:scale-dependence}

Taken together, the results of this chapter show that ``the registration is
valid'' is not one claim but four, holding at different scales and by
different evidence:

\begin{center}
\small
\begin{tabular}{@{}p{0.19\linewidth}p{0.36\linewidth}p{0.33\linewidth}@{}}
\toprule
\textbf{Level} & \textbf{Question} & \textbf{Evidence in this report} \\
\midrule
Surface correspondence & Are the two surfaces geometrically close? &
  Chamfer, F-score (misleading alone, \F{F1}) \\
Anatomical correspondence & Does model location $i$ mark the same structure
  across specimens? & Synthetic oracle, corrected expert benchmark
  (\F{F3}, \F{F6}) \\
Parametric correspondence & Do recovered $\bm\beta,\bm\theta$ mean the same
  thing across specimens? & Only partially established: 1 of 49
  joint-derived features reaches $r=0.5$ \\
Measurement validity & Does a derived measurement match an independent one? &
  Atta head width, allometric slopes (\S\ref{sec:measurement}) \\
\bottomrule
\end{tabular}
\end{center}

A model can therefore give accurate population-level and head-width
measurements while failing to preserve the individual-level correspondence
needed for specimen-by-specimen comparison of distal segment lengths. Both
statements are true, and each is supported only at its own level.
